% answers_short.tex —— 2006 年·普通高等学校招生考试·广东卷 · 简短答案速查
% 由 tools/make_short_answers.py 自 answers.tex 的 \gksolhead/\gkanswer 提取，
% 并与 2006/2006_数学真题_广东_答案要点.md 逐题比对后生成，勿手改（改源 MD 后重跑）。

\begin{gknotice}
\textbf{使用说明}：本节只给最终答案便于快速核对，完整推导见解析版。答案为编写者独立推导，\textbf{非官方标准答案}。
\end{gknotice}

\begin{gknotice}
选择题答案中的 U 表示待核，未给出猜测答案，详情见解析版。
\end{gknotice}

\gkpart{选择题}
\begin{center}\begin{tabular}{ccccccccccc}
\toprule
题号 & 1 & 2 & 3 & 4 & 5 & 6 & 7 & 8 & 9 & 10 \\
\midrule
答案 & B & D & A & A & B & U & C & C & D & B \\
\bottomrule
\end{tabular}\end{center}

\gkpart{填空题}
\begin{itemize}[leftmargin=2.6em]
  \item[\textbf{11.}] $1/4$。
  \item[\textbf{12.}] $27\pi$。
  \item[\textbf{13.}] $-1320$。
  \item[\textbf{14.}] $f(3)=10$，$f(n)=\binom{n+2}{3}$。
\end{itemize}

\gkpart{解答题}
\begin{itemize}[leftmargin=2.6em]
  \item[\textbf{15.}] (1) $2\pi$；(2) 最大值 $\sqrt2$，最小值 $-\sqrt2$；(3) $-7/16$
  \item[\textbf{16.}] (1) U/PENDING；若补充假设两次射击独立，则为 $0.04$；(2) U/PENDING；题面没有定义随机变量 $\xi$；(3) U/PENDING；题面没有定义随机变量 $\xi$
  \item[\textbf{17.}] (1) $45^\circ$；(2) $\arccos\left(\dfrac{\sqrt{82}}{10}\right)$
  \item[\textbf{18.}] (1) $A=(-1,0),\ B=(1,4)$；(2) $(x-8)^2+(y+2)^2=9$
  \item[\textbf{19.}] (1) U/PENDING；总和与偶数项子列和矛盾；(2) U/PENDING；依赖不存在的数列参数；(3) U/PENDING；依赖不存在的数列参数
  \item[\textbf{20.}] (1) U/PENDING；题面给定区间使 $\varphi(2x)$ 超出其定义域；(2) 唯一；(3) U/PENDING；递推方向与任意给定初值不相容，见解析
\end{itemize}
